\pagenumbering{arabic}
\chapter{Introduction}

Organic electronic materials have attracted significant attention in recent years due to their potential applications. The organic electronics include, e.g., \acp{OLED}, \acp{OPV}, and \acp{OFET}.\autocite{Mitschke2000,Farinola2011,Dou2015,Wu2010,Lu2015,Koch2007}

Generally, an extended $\pi$-conjugated system is considered to be a prerequisite for high electron mobility and thus for good performance in organic electronics.\autocite{Jiang2022,Zhang2018} Model molecules with an extended $\pi$-conjugated system, such as pentacene, have been widely studied for its application in organic electronic materials.\autocite{Kitamura2008,Yang2016,Kitamura2003}

A structurally analogous molecule is \acf{QA}, which possesses a comparable molecular structure but contains four distinct functional groups. \Ac{QA} is a planar, prochiral organic molecule (\ce{C20H12N2O2}), which is widely used as pigment (Violet 19) in paints, laquers, plastics, and printing inks.\autocite{Struve1958,Zollinger2003, Faulkner2009,Hye2012}

However, \ac{QA} is not only an interesting molecule for its application as pigment, but also for its potential application in organic electronics. The functional groups of \ac{QA} are responsible for the intermolecular hydrogen bonds, which lead to high field effect mobilities and photocurrents.\autocite{Glowacki2011,Glowacki2012,Glowacki2013} More detailed, \ac{QA} has been studied in the context of \acp{OLED}\autocite{Wang2016a,Min2021,Cunha2018}, \acp{OFET}\autocite{Jeong2018,Jeong2017}, and \acp{OPV}\autocite{Dunst2017,Sung2017,Kanbur2019,Berg2009}.
 
In addition, \ac{QA} has been shown to form self-ordering molecular layers on metal surfaces, such as the Ag(100), the Ag(111), and the Cu(111) surface.\autocite{Humberg2020,Humberg2024,Humberg2025,Wagner2014} A common structural motif of \ac{QA} on these surfaces is the formation of 1D molecular chains, which is also observed in the bulk structure of \ac{QA}.\autocite{Paulus2007} The formation of these molecular chains is driven by intermolecular hydrogen bonds between the amine and carbonyl groups.\autocite{Humberg2020,Humberg2024,Humberg2025}

The investigation of \ac{QA} on the Ag(100) surface revealed two different phases: the $\alpha$-phase and the $\beta$-phase. The metastable $\alpha$-phase with a \ac{pol} structure consists of the previously mentioned homochiral 1D molecular chains, which are parallel to each other. The thermodynamically favored $\beta$-phase with a commensurate structure is formed irreversibly by annealing the sample and consists of dimers with alternating handedness. In addition, the $\beta$-phase is supposed to have fewer hydrogen bonds than the $\alpha$-phase, which is energetically unfavorable, but is compensated by a stronger molecule-substrate interaction.\autocite{Humberg2020,Humberg2025}

Nevertheless, the phase transition, including the differences between the two phases and the hydrogen bonding, is not fully understood and requires further investigation. This is due to the fact that the previous studies of \ac{QA} on the Ag(100) surface were mainly based on \ac{STM} and \ac{LEED} measurements, which only provide structural information.\autocite{Humberg2020,Humberg2025}

Consequently, this thesis focuses on how the two phases of \ac{QA} on the Ag(100) surface is affected by the interplay between intemolecular interactions via hydrogen bonds and adsorbate-susbtrate interactions. Especially, the formation of the hydrogen bonds and their change between the two phases, as well as the specific atoms that are involved in the adsorbate-substrate interaction are of interest.

To investigate these questions, \ac{XPS} and \ac{NIXSW} techniques are used. \Ac{XPS} provides insights into the binding energies of the core level electrons, which contain information about the chemical environment of the different atoms in the molecule. Furthermore, \ac{NIXSW} allows to determine the adsorption height with high precision, which is directly related to the strength of the adsorbate-substrate interaction.

As a result thereof, new insights into the situation of the hydrogen bonds are obtained. Previous studies proposed two hydrogen bonds per molecule in the $\alpha$-phase and 1.5 hydrogen bonds per molecule in the $\beta$-phase.\autocite{Humberg2020,Humberg2025} The \ac{XPS} support the two hydrogen bonds per molecule in the $\alpha$-phase, but contradict the 1.5 hydrogen bonds per molecule in the $\beta$-phase, as the results indicate only one hydrogen bonds per molecule. In combination with the adsorption heights obtained with \ac{NIXSW}, which show strong changes for the N atoms, a new structure model for \ac{QA} on the Ag(100) surface is necessary.

\cleardoublepage